\documentclass[10pt,a4paper]{amsart}
\usepackage[margin=19mm]{geometry}
\usepackage{amsmath,amssymb,mathtools}
\usepackage[scr=rsfs,cal=euler]{mathalfa}
\usepackage{xfrac}
\usepackage[unicode,hidelinks,bookmarks=false]{hyperref}
\newcommand{\ie}{i.e.\ }
\newcommand{\cf}{cf.\ }
\newcommand{\resp}{respectively\ }
\newcommand{\Subsection}{\subsection}
\providecommand{\nobreakdash}{\nobreak\hbox{-}}
\setlength{\emergencystretch}{2em}
\multlinegap0pt
\usepackage{amssymb}
\usepackage{enumerate}
\newtheorem{lemma}{Lemma}
\title{On divisor sums due to Erd\H{o}s and Ramanujan}
\author{John M. Campbell}
\date{Source: arXiv:2605.00695v1 (CC BY 4.0)}
\begin{document}
\maketitle

\noindent\emph{Source and licence.} On divisor sums due to Erd\H{o}s and Ramanujan. Version arXiv:2605.00695v1 (CC BY 4.0), distributed by arXiv under the Creative Commons Attribution 4.0 International licence, http://creativecommons.org/licenses/by/4.0/, as recorded in the arXiv licence metadata for this exact version and shown on the abstract page https://arxiv.org/abs/2605.00695v1. Full licence notice: \texttt{licenses/arxiv-2605.00695-CC-BY-4.0.txt}.\par\smallskip

\noindent\textit{Editorial scope.} Counted unit: the article's lemma LemmaLower (source0.tex lines 359--361). The article's complete original proof of it is delivered with it (source0.tex lines 363--383, one \texttt{proof} environment of 21 lines). No mathematics is written, repaired or re-derived: the statement and the proof are the article's own lines, in the article's own notation, and no retained line is substituted or re-flowed.  Retained uncounted as the source-local context the proof needs: lines 283--286.  Nothing was omitted from any retained span, so the package carries no omission record.  No retained line contains a citation, so the wrapper carries no bibliography and no bibliography entry.  No retained line contains a figure environment, an image-inclusion command or drawing code, so no image is delivered and the package declares no asset dependency.  Editorial formatting only, in addition: an AMS \texttt{amsart} preamble that restates the article's own declarations and macros used by the retained lines, namely the environments \texttt{lemma} on separate counters, together with the standard packages those lines need.  The retained environments print in this document as Lemma 1 in order of appearance, whereas the article prints the same items with the numbering of its own full document.  The complete native source of the article as distributed in the arXiv e-print for this exact version is bundled unchanged as \texttt{sources/199584/source0.tex}.
\begin{lemma}\label{LemmaB2}
Let $$ \mathcal{B}_2(x) := \left\{n \leq x:\omega(n) \notin \left[\frac12\log\log x,2\log\log x\right]\right\}. $$ Then $$ \sum_{n \in 
 \mathcal{B}_2(x)} \frac{1}{d(d(n))} \ll \frac{x}{\log\log x}. $$
\end{lemma}

\begin{lemma}\label{LemmaLower}
 We have that $$ \frac{x}{\log\log x} \ll \sum_{n\le x}\frac{1}{d(d(n))}. $$
\end{lemma}

\begin{proof}
Let
\begin{equation}\label{displayS2}
\mathcal{B}_{3}(x) := \left\{ n \le x : n \text{ squarefree},\ 
\omega(n) \in \left[\frac12\log\log x, 2\log\log x\right] \right\}.
\end{equation}
 It is classically and well known that $ \#\{n \le x : n \text{ squarefree}\} = \frac{6}{\pi^2}x + O(\sqrt{x})$. 
 By rewriting \eqref{displayS2} as a set difference involving $\mathcal{B}_{2}$, we find that
\begin{equation}\label{fromsetdiff}
 \#\mathcal{B}_{3}(x) \geq \#\{ n \leq x : \text{$n$ squarefree} \} 
 - \#\mathcal{B}_{2}(x). 
\end{equation}
 It was shown in the proof of Lemma \ref{LemmaB2} that $ \#\mathcal{B}_{2}(x) \ll \frac{x}{\log \log x}$. This together with the classical 
 estimate for $\#\{n \le x : n \text{ squarefree}\}$ then give us, from \eqref{fromsetdiff}, that $\#\mathcal{B}_{3}(x) \gg x$. Now, let 
 $n \in \mathcal{B}_{3}(x)$. Since \(n\) is squarefree, we have that the relation $ d(n) = 2^{\omega(n)}$ holds. As a consequence, we 
 find that the equalities such that $ d(d(n)) = d\big(2^{\omega(n)}\big)=\omega(n)+1 $ hold true. Moreover, the definition of 
 $\mathcal{B}_{3}(x)$ gives us that $ \omega(n)+1\leq 2\log\log x+1 \ll \log\log x$, so that $ \frac{1}{d(d(n))} = \frac{1}{\omega(n) + 
 1} \gg \frac{1}{\log\log x}$. It then follows that $$ \sum_{n \leq x} \frac{1}{d(d(n))} \geq \sum_{n\in\mathcal{B}_{3}(x)}\frac{1}{d(d(n))}
\gg \frac{\#\mathcal{B}_{3}(x)}{\log\log x} \gg
\frac{x}{\log\log x}, $$ giving us the desired result. 
\end{proof}

\end{document}
